\documentclass[11pt,a4paper]{article}
\usepackage{gtu-el-yazisi}
\begin{document}
\notbilgi{Görüntü İşleme}{Ders 12: Granulometry, geometrik dönüşüm, feature extraction}{Ders-12.pdf}

\sayfa{1}
\section{\okunamadi{}}
Parçacık sayma \hfill Granulometry

\elcizimi[0.35]{sekiller/Ders-12-s1-1.png}[1]

opening erosion, dialation

\sekilnotu[1]{Bu sayfa basılı materyal içeriyor (ör. ders slaytı); dijitalleştirilmedi.}

yapı elemanı her adımda büyütülür \\
Her pikselin değeri toplanır \\
birbirinden çıkarırsak

ne büyüklükte objeler var

\okunamadi{Router chef} segmentasyon

\sayfa{2}
\section{Geometric Transformasyon}
\elcizimi[0.4]{sekiller/Ders-12-s2-1.png}[2]

hepsini aynı boyutta görmez

\section{Geometric Transformation}
\elcizimi[0.6]{sekiller/Ders-12-s2-2.png}[2]

\elcizimi[0.6]{sekiller/Ders-12-s2-3.png}[2]

pseudo-code

for $x=0$; nrows in target \\
\hspace*{1.5em} for $y=0$; ncols in target \\
\hspace*{3em} compute $\hat{x}$ = inverse spatial transformation of $x$ \\
\hspace*{3em} $\hat{y}$ = inverse spatial Transformation $y$ \\
\hspace*{3em} $f(x,y)$ = gray level interpolation of $g(\hat{x},\hat{y})$

\sayfa{3}
\underline{scaling}: \quad $\hat{x}=\dfrac{x}{c}$ \qquad\quad $\hat{y}=\dfrac{y}{c}$ \hfill {\small $c$ büyünce görüntü büyür}

\underline{rotation}:
\begin{align*}
\hat{x} &= x\cos\theta - y\sin\theta\\
\hat{y} &= x\sin\theta + y\cos\theta
\end{align*}
\[
\begin{bmatrix}\hat{x}\\ \hat{y}\end{bmatrix}
= \begin{bmatrix}\cos\theta & -\sin\theta\\ \sin\theta & \cos\theta\end{bmatrix}
\begin{bmatrix}x\\ y\end{bmatrix}
\]

\section{Gray level interpolasyon}
level 0: nearest neighbour (NN) \\
level 1: bilinear interpolation (BI)

\begin{align*}
\hat{x} &= k + r_x & \qquad \hat{x} &= \tfrac{4}{3} = \overset{k}{1} + \overset{r_x}{\tfrac{1}{3}}\\
\hat{y} &= j + r_y & \qquad \hat{y} &= \tfrac{15}{2}\footnotemark = \underset{k}{2} + \underset{r_y=3/6}{\tfrac{3}{6}}
\end{align*}
\footnotetext{Claude notu: muhtemelen $\tfrac{15}{6}$ olmalı.}

\elcizimi[0.35]{sekiller/Ders-12-s3-1.png}[3]

\elcizimi[0.5]{sekiller/Ders-12-s3-2.png}[3]

\elcizimi[0.4]{sekiller/Ders-12-s3-3.png}[3]

\sayfa{4}
\section{Örnek}
\noindent
\begin{minipage}[t]{0.48\linewidth}
\begin{tabular}{c|cccc}
 & 0 & 1 & 2 & 3 \\ \hline
0 & 1 & 1 & 100 & 1 \\
1 & 1 & 1 & 100 & 1 \\
2 & 100 & 100 & 100 & 100 \\
3 & 1 & 1 & 100 & 1
\end{tabular}

\bigskip
\begin{tabular}{c|ccc}
 & \multicolumn{3}{c}{$y$} \\
 & 0 & 1 & 2 \\ \hline
0 & & & \\
$x$\ 1 & & & \\
2 & & &
\end{tabular}
\end{minipage}\hfill
\begin{minipage}[t]{0.48\linewidth}
Sonuç

\begin{tabular}{c|ccc}
 & 0 & 1 & 2 \\ \hline
0 & 1 & 34 & 34 \\
$x$\ 1 & 34 & 56 & 56 \\
2 & 34 & 56 & 56
\end{tabular}
\end{minipage}

\medskip
$c=\dfrac{3}{4}$ \qquad $c=\dfrac{2}{3}$ alma

\medskip
\begin{center}
\begin{tabular}{cc|cc|c|c}
$x$ & $y$ & $\hat{x}$ & $\hat{y}$ & NN & BI \\ \hline
0 & 0 & 0 & 0 & {\scriptsize 0,0}\ 1 & \\
0 & 1 & 0 & $4/3$ & {\scriptsize 0,1}\ 1 & \\
0 & 2 & 0 & $8/3$ & {\scriptsize 0,3}\ 1 & \\
1 & 0 & $4/3$ & 0 & 1 & \\
1 & 1 & $4/3$ & $4/3$ & 1 & \\
1 & 2 & $4/3$ & $8/3$ & 1 & \\
2 & 0 & $8/3$ & 0 & 1 & \\
2 & 1 & $8/3$ & $4/3$ & 1 & \\
2 & 2 & $8/3$ & $8/3$ & 1 & \\
\end{tabular}
\end{center}

$x=1$ \qquad $y=2$ \\
$\hat{x}=\dfrac{4}{3}$ \qquad $\hat{y}=\dfrac{8}{3}$ \\
$\hat{x}=1+\dfrac{1}{3}$ \qquad $\hat{y}=2+\dfrac{2}{3}$

\elcizimi[0.45]{sekiller/Ders-12-s4-1.png}[4]

\sayfa{5}
\elcizimi[0.4]{sekiller/Ders-12-s5-1.png}[5]

\section{Feature Extraction}
\elcizimi[0.5]{sekiller/Ders-12-s5-2.png}[5]

Moment hesabı yapabiliriz

fourier discriptor

\elcizimi[0.3]{sekiller/Ders-12-s5-3.png}[5]

\[ s(r) = x + jy \]
\[ S[k] = \sum_{r=0}^{R} s(r)\, e^{-j2\pi k/R} \]

Texture analizi \qquad\qquad Co-occurance Matrix

\[ P_{\phi,d}(a,b) \]

\elcizimi[0.55]{sekiller/Ders-12-s5-4.png}[5]

\elcizimi[0.4]{sekiller/Ders-12-s5-5.png}[5]

\end{document}
